\frfilename{mt134.tex}
\versiondate{7.1.04}
\copyrightdate{1994}

%B(x,\delta) not B(x;\delta)
\def\chaptername{Pelengkap}
\def\sectionname{Lebih lanjut tentang ukuran Lebesgue}

\newsection{134}

Sifat-sifat khusus ukuran Lebesgue akan mengisi bagian yang cukup besar dari
risalah ini.
Dalam bagian ini saya menyajikan aneka hasil dasar yang relatif mudah. Dalam
134A-134F, %134A 134B 134D 134E 134F
$r$ akan merupakan suatu bilangan bulat tetap yang lebih besar daripada atau
sama dengan $1$, $\mu$ akan merupakan ukuran Lebesgue pada $\BbbR^r$, dan
$\mu^*$ akan merupakan ukuran luar Lebesgue\cmmnt{ (lihat 132C)}; ketika saya
mengatakan bahwa suatu himpunan atau fungsi `terukur', hendaknya dipahami
bahwa (kecuali dinyatakan lain) artinya `terukur terhadap aljabar-$\sigma$
himpunan-himpunan terukur Lebesgue', sedangkan `terabaikan' berarti
`terabaikan untuk ukuran Lebesgue'. Sebagian besar hasil akan dinyatakan
dalam istilah yang disesuaikan dengan kasus multidimensi; tetapi jika minat
utama Anda adalah garis bilangan real, Anda tidak akan kehilangan satu pun
gagasan apabila membaca seluruh bagian ini seolah-olah $r=1$.

\leader{134A}{Proposisi} Baik ukuran luar Lebesgue maupun ukuran Lebesgue
invarian terhadap translasi; yaitu, dengan menetapkan
$A+x=\{a+x:a\in A\}$ untuk $A\subseteq\BbbR^r$, $x\in\BbbR^r$, kita
memiliki

(a) $\mu^*(A+x)=\mu^* A$ untuk setiap $A\subseteq\BbbR^r$,
$x\in\BbbR^r$;

(b) jika $E\subseteq\BbbR^r$ terukur dan $x\in\BbbR^r$, maka $E+x$
terukur, dengan $\mu(E+x)=\mu E$.

\proof{ Intinya adalah bahwa jika $I\subseteq\BbbR^r$ merupakan interval
setengah terbuka, sebagaimana didefinisikan dalam 114Aa/115Ab, maka $I+x$
juga demikian, dan $\lambda(I+x)=\lambda I$ untuk setiap $x\in\BbbR^r$,
dengan $\lambda$ didefinisikan seperti dalam 114Ab/115Ac; hal ini langsung
mengikuti definisi, karena
$\coint{a,b}+x=\coint{a+x,b+x}$.

\medskip

{\bf (a)} Jika $A\subseteq\BbbR^r$, $x\in\BbbR^r$, dan $\epsilon>0$,
kita dapat menemukan barisan $\langle I_j\rangle_{j\in\Bbb N}$ interval
setengah terbuka sedemikian sehingga $A\subseteq\bigcup_{j\in\Bbb N}I_j$
dan $\sum_{j=0}^{\infty}\lambda I_j\le\mu^* A+\epsilon$. Sekarang
$A+x\subseteq\bigcup_{j\in\Bbb N}(I_j+x)$, sehingga


\Centerline{$\mu^*(A+x)\le\sum_{j=0}^{\infty}\lambda
(I_j+x)=\sum_{j=0}^{\infty}\lambda I_j\le\mu^* A+\epsilon$.}

\noindent Karena $\epsilon$ sembarang,
$\mu^* (A+x)\le\mu^* A$. Demikian pula,

\Centerline{$\mu^*A=\mu^*((A+x)+(-x))\le\mu^*(A+x)$,}

\noindent sehingga $\mu^*(A+x)=\mu^* A$, seperti yang dinyatakan.

\wheader{134A}{6}{2}{2}{12pt}

{\bf (b)} Sekarang andaikan $E\subseteq\BbbR^r$ terukur,
$x\in\BbbR^r$, dan $A\subseteq\BbbR^r$. Maka, dengan menggunakan (a)
berulang kali,

$$\eqalign{\mu^*(A\cap(E+x))+\mu^*(A\setminus(E+x))
&=\mu^*(((A-x)\cap E)+x)+\mu^*(((A-x)\setminus E)+x)\cr
&=\mu^*((A-x)\cap E)+\mu^*((A-x)\setminus E)\cr
&=\mu^*(A-x)
=\mu^* A,\cr}$$

\noindent dengan menuliskan $A-x$ untuk
$A+(-x)=\{a-x:a\in A\}$. Karena $A$ sembarang, $E+x$ terukur. Sekarang

\Centerline{$\mu(E+x)=\mu^*(E+x)=\mu^* E=\mu
E$.}
}%end of proof of 134A



\leader{134B}{Teorema} Tidak setiap subhimpunan dari $\BbbR^r$ terukur
Lebesgue.

\proof{ Tetapkan $\tbf{0}=(0,\ldots,0)$,
$\tbf{1}=(1,\ldots,1)\in\BbbR^r$. Pada

\Centerline{$\coint{\tbf{0},\tbf{1}}
=\{(\xi_1,\ldots,\xi_r):\xi_i\in\coint{0,1}$ untuk setiap $i\le r\}$,}

\noindent tinjau relasi $\sim$, yang didefinisikan dengan mengatakan bahwa
$x\sim y$ jika dan hanya jika $y-x\in\BbbQ^r$. Mudah dilihat bahwa ini
merupakan relasi ekuivalensi, sehingga membagi
$\coint{\tbf{0},\tbf{1}}$ menjadi kelas-kelas ekuivalensi. Pilih satu
titik dari setiap kelas ekuivalensi ini, dan misalkan $A$ adalah himpunan
titik-titik yang diperoleh dengan cara tersebut. Maka
$\mu^* A\le\mu^*\coint{\tbf{0},\tbf{1}}=1$.

Tinjau $A+\BbbQ^r=\{a+q:a\in A,\,q\in\BbbQ^r\}
=\bigcup_{q\in\BbbQ^r}A+q$. Himpunan ini sama dengan $\BbbR^r$.
\Prf\ Jika $x\in\BbbR^r$, terdapat $e\in\BbbZ^r$ sedemikian sehingga
$x-e\in\coint{\tbf{0},\tbf{1}}$; terdapat $a\in A$ sedemikian sehingga
$a\sim x-e$, yaitu $x-e-a\in\BbbQ^r$; sekarang
$x=a+(e+x-e-a)\in A+\BbbQ^r$.\ \Qed\ Selanjutnya, $\BbbQ^r$
terhitung (111F(b-iv)), sehingga kita memiliki

\Centerline{$\infty=\mu \BbbR^r\le\sum_{q\in\BbbQ^r}\mu^*(A+q)$,}

\noindent dan pasti terdapat suatu $q\in\BbbQ^r$ sedemikian sehingga
$\mu^*(A+q)>0$; tetapi karena $\mu^*$ invarian terhadap translasi (134A),
$\mu^* A>0$.

Ambil $n\in\Bbb N$ sedemikian sehingga $n>2^r/\mu^* A$, dan ambil
$q_1,\ldots,q_n\in\coint{\tbf{0},\tbf{1}}\cap\BbbQ^r$ yang berbeda-beda.
Jika $a$, $b\in A$ dan $1\le i<j\le n$, maka
$a+q_i\ne b+q_j$; sebab jika $a=b$, maka $q_i\ne q_j$, sedangkan jika
$a\ne b$, maka $a\not\sim b$, sehingga $b-a\ne q_i-q_j$. Dengan demikian
$A+q_1,\ldots,A+q_n$ saling lepas. Di sisi lain, semuanya merupakan
subhimpunan $\coint{\tbf{0},\tbf{2}}$. Jadi kita memiliki

\Centerline{$\sum_{i=1}^n\mu^*(A+q_i)
=n\mu^* A>2^r=\mu\coint{\tbf{0},\tbf{2}}\ge\mu^*(\bigcup_{1\le i\le
n}(A+q_i))$.}

\noindent Akibatnya, tidak semua $A+q_i$ dapat terukur; karena ukuran
Lebesgue invarian terhadap translasi, kita melihat bahwa $A$ sendiri tidak
terukur. Bagaimanapun, kita telah menemukan suatu himpunan tak-terukur.
}%end of proof of 134B


\cmmnt{
\leader{*134C}{Catatan} 134B dikenal sebagai `konstruksi Vitali'.

Perhatikan bahwa pada awal pembuktian saya meminta Anda {\it memilih} satu
anggota dari setiap kelas ekuivalensi untuk $\sim$. Tentu saja ini merupakan
penggunaan Aksioma Pilihan. Sejauh ini saya agak jarang menggunakan aksioma
pilihan. Salah satunya terdapat dalam (a-iv) pada pembuktian 114D/115D;
penggunaan yang lebih awal terdapat dalam 112Db; dan satu lagi dalam 121A.
Lihat juga 1A1F. Dalam semua kasus tersebut, yang terlibat hanya `pilihan
terhitung'; yaitu, saya perlu memilih secara serentak satu anggota dari setiap
himpunan dalam suatu barisan yang telah ditentukan. Karena pasti terdapat tak
terhitung banyak kelas ekuivalensi untuk $\sim$, bentuk pilihan yang diperlukan
untuk contoh di atas pada hakikatnya lebih kuat daripada yang diperlukan untuk
hasil-hasil positif sejauh ini. Bahkan, bagian yang sangat besar dari teori
ukuran dapat dikembangkan tanpa menggunakan kekuatan penuh aksioma pilihan.

Arti penting hal ini ialah bahwa mungkin terdapat suatu sistem matematika
konsisten yang mengakui bentuk Aksioma Pilihan yang cukup kuat untuk
memungkinkan teori ukuran, tetapi tidak cukup kuat untuk membangun suatu
himpunan tak-terukur Lebesgue. Sistem semacam itu memang telah dikembangkan oleh
R.M.Solovay ({\smc Solovay 70}). (Dalam arti formal masih terdapat ruang untuk
keraguan tersisa mengenai konsistensinya. Menurut saya, ini tidak penting.)
Dalam Jilid 5 saya akan kembali ke pertanyaan mengenai seperti apa ukuran
Lebesgue dengan aksioma pilihan yang lemah, atau tanpa aksioma pilihan sama
sekali. Untuk saat ini, saya harus mengatakan bahwa hampir seluruh teori
ukuran terus berjalan dalam arah yang setidaknya konsisten dengan aksioma
pilihan penuh, sehingga himpunan-himpunan tak-terukur senantiasa hadir,
setidaknya secara potensial; dan itulah sikap normal saya dalam risalah ini.
Namun saya menyebutkan pokok ini pada tahap awal karena saya percaya bahwa
sewaktu-waktu pusat perhatian dapat beralih ke sistem-sistem tempat aksioma
pilihan tidak berlaku; dalam hal itu teori ukuran tanpa himpunan tak-terukur
dapat menjadi penting bagi banyak matematikawan murni, bahkan matematikawan
terapan, yang tidak mempunyai alasan untuk setia pada aksioma pilihan selain
kemudahan dapat mengutip hasil dari buku-buku seperti ini.

Perlu saya catat bahwa meskipun kita memerlukan bentuk aksioma pilihan yang
cukup kuat untuk membangun himpunan tak-terukur Lebesgue, suatu himpunan
non-Borel dapat dibangun dalam teori-teori himpunan yang jauh lebih lemah.
Salah satu konstruksi yang mungkin diuraikan dalam \S423 di Jilid 4.

Tentu saja terdapat subhimpunan tak-terukur Lebesgue dari $\Bbb R$ jika dan
hanya jika terdapat fungsi tak-terukur Lebesgue dari $\Bbb R$ ke $\Bbb R$;
sebab jika setiap himpunan terukur, definisi 121C memperjelas bahwa setiap
fungsi bernilai real pada sembarang subhimpunan dari $\Bbb R$ terukur; sedangkan
jika $A\subseteq\Bbb R$ tidak terukur, maka
$\chi A:\Bbb R\to\Bbb R$ tidak terukur.
}%end of comment

\leader{*134D}{}\cmmnt{ Sebenarnya terdapat hasil-hasil yang jauh lebih kuat
daripada 134B mengenai keberadaan himpunan tak-terukur (tentu saja, asalkan
kita memperbolehkan diri menggunakan aksioma pilihan). Di sini saya memberikan
satu hasil yang dapat dicapai melalui sedikit penyempurnaan metode-metode 134B.

\medskip

\noindent}{\bf Proposisi} Terdapat himpunan $C\subseteq\BbbR^r$ sedemikian
sehingga $F\cap C$ tidak terukur untuk setiap himpunan terukur $F$ yang
ukurannya tak nol; dengan demikian, baik $C$ maupun komplemennya berukuran
luar penuh dalam $\BbbR^r$.

\proof{{\bf (a)} Mulailah dengan suatu himpunan
$A\subseteq\coint{\tbf{0},\tbf{1}}\subseteq\BbbR^r$ sedemikian sehingga
$\langle A+q\rangle_{q\in\BbbQ^r}$ merupakan partisi atas $\BbbR^r$,
sebagaimana dibangun dalam pembuktian 134B. Seperti dalam 134B, ukuran luar
$\mu^*A$ dari $A$ harus lebih besar daripada $0$. Argumen di sana sebenarnya
menunjukkan bahwa $\mu F=0$ untuk setiap himpunan terukur $F\subseteq A$.
\Prf\ Untuk setiap $n$ kita dapat menemukan
$q_1,\ldots,q_n\in\coint{\tbf{0},\tbf{1}}\cap\BbbQ^r$ yang berbeda-beda,
dan sekarang

\Centerline{$n\mu F=\mu(\bigcup_{1\le i\le
n}F+q_i)\le\mu\coint{\tbf{0},\tbf{2}}=2^r$,}

\noindent sehingga $\mu F\le 2^r/n$; karena $n$ sembarang, $\mu F=0$.
\Qed

\medskip

{\bf (b)} Sekarang misalkan $E\subseteq\coint{\tbf{0},\tbf{1}}$
merupakan selubung terukur dari $A$ (132Ef). Maka $E+q$ merupakan selubung
terukur dari $A+q$ untuk setiap $q$. \Prf\ Saya harap hal ini akan segera
menjadi `konsekuensi jelas dari invariansi ukuran Lebesgue terhadap translasi'.
Secara terperinci: $A+q\subseteq E+q$, $E+q$ terukur dan, untuk setiap $F$
terukur,

$$\eqalign{\mu(F\cap(E+q))
&=\mu(((F-q)\cap E)+q)
=\mu((F-q)\cap E)\cr
&=\mu^*((F-q)\cap A)
=\mu^*(((F-q)\cap A)+q)
=\mu^*(F\cap(A+q)),\cr}$$

\noindent dengan menggunakan 134A berulang kali.\ \Qed\ Selain itu, $E$
merupakan selubung terukur dari $A'=E\setminus A$.
\Prf\ Tentu saja $E$ merupakan himpunan terukur yang memuat $A'$. Jika
$F\subseteq E\setminus A'$ terukur, maka $F\subseteq A$, sehingga
$\mu F=0$, menurut (a); sekarang 132Ea memberi tahu kita bahwa $E$ merupakan
selubung terukur dari $A'$.\ \Qed\ Akibatnya, $E+q$ merupakan selubung
terukur dari $A'+q$ untuk setiap $q$.

\medskip

{\bf (c)} Misalkan $\sequencen{q_n}$ suatu barisan yang menjangkau
$\BbbQ^r$. Maka

\Centerline{$\bigcup_{n\in\Bbb N}E+q_n
\supseteq\bigcup_{n\in\Bbb N}A+q_n=\BbbR^r$.}

\noindent Tuliskan $E_n$ untuk
$E+q_n\setminus\bigcup_{i<n}E+q_i$ bagi $n\in\Bbb N$, sehingga
$\sequencen{E_n}$ saling lepas dan
$\bigcup_{n\in\Bbb N}E_n=\BbbR^r$.

Sekarang tetapkan

\Centerline{$C=\bigcup_{n\in\Bbb N}E_n\cap(A+q_n)$.}

\noindent Inilah himpunan dengan sifat-sifat yang diperlukan.

\Prf\ {\bf (i)} Misalkan $F\subseteq\BbbR^r$ sembarang himpunan terukur
tak-terabaikan. Maka pasti terdapat suatu $n\in\Bbb N$ sedemikian sehingga
$\mu(F\cap E_n)>0$. Namun ini berarti bahwa

\Centerline{$\mu^*(F\cap E_n\cap C)\ge\mu^*(F\cap E_n\cap(A+q_n))
=\mu(F\cap E_n\cap(E+q_n))=\mu(F\cap E_n)$,}

$$\eqalign{\mu^*(F\cap E_n\setminus C)
&\ge\mu^*(F\cap E_n\cap((E+q_n)\setminus(A+q_n)))\cr
&=\mu(F\cap E_n\cap(E+q_n))
=\mu(F\cap E_n).\cr}$$

\noindent Karena

\Centerline{$\mu(F\cap E_n)\le\mu(E+q_n)=\mu E\le 1$,}

\noindent
$\mu^*(F\cap E_n\cap C)+\mu^*(F\cap E_n\setminus C)>\mu(F\cap E_n)$,
dan $F\cap C$ tidak mungkin terukur.

\medskip

\quad{\bf (ii)} Khususnya, tidak ada subhimpunan terukur dari
$\BbbR^r\setminus C$ yang dapat memiliki ukuran tak nol, dan $C$ berukuran
luar penuh; demikian pula, $C$ tidak memiliki subhimpunan terukur berukuran
tak nol, dan $\BbbR^r\setminus C$ berukuran luar penuh.\ \Qed
}%end of proof of 134D

\cmmnt{\medskip

\noindent{\bf Catatan} Sebenarnya, untuk setiap barisan
$\sequencen{D_n}$ yang terdiri atas subhimpunan-subhimpunan dari $\BbbR^r$,
terdapat himpunan
$C\subseteq\BbbR^r$ sedemikian sehingga

\Centerline{$\mu^*(E\cap D_n\cap C)=\mu^*(E\cap D_n\setminus C)
=\mu^*(E\cap D_n)$}

\noindent untuk setiap himpunan terukur $E\subseteq\BbbR^r$ dan setiap
$n\in\Bbb N$. Namun pembuktian hasil ini harus menunggu Jilid 5.
}%end of comment

\cmmnt{
\leader{134E}{Himpunan Borel dan ukuran Lebesgue pada $\BbbR^r$}
Ingat dari 111G bahwa keluarga $\Cal B$ himpunan-himpunan Borel dalam
$\BbbR^r$ adalah aljabar-$\sigma$ yang dibangkitkan oleh keluarga
himpunan-himpunan terbuka. Dalam 114G/115G saya menunjukkan bahwa setiap
himpunan Borel dalam $\BbbR^r$ terukur Lebesgue. Sudah waktunya kita kembali
ke topik ini dan melihat lebih dekat hubungan yang sangat erat antara
himpunan Borel dan himpunan terukur.

Ingat bahwa suatu himpunan $A\subseteq\BbbR^r$ disebut {\bf berbatas} jika
terdapat suatu $M$ sedemikian sehingga
$A\subseteq B(\tbf{0},M)=\{x:\|x\|\le M\}$; secara ekuivalen, jika
$\sup_{x\in A}|\xi_j|<\infty$ untuk setiap $j\le r$ (dengan menuliskan
$x=(\xi_1,\ldots,\xi_r)$, seperti dalam \S115).
}%end of comment

\leader{134F}{Proposisi} (a) Jika $A\subseteq\BbbR^r$ sembarang
himpunan, maka

\Centerline{$\mu^*A
=\inf\{\mu G:G\text{ terbuka},\,G\supseteq A\}
=\min\{\mu H:H\text{ Borel},\,H\supseteq A\}$.}

(b) Jika $E\subseteq\BbbR^r$ terukur, maka

\Centerline{$\mu E=\sup\{\mu F:F
\text{ tertutup dan berbatas},\,F\subseteq E\}$,}

\noindent dan terdapat himpunan-himpunan Borel $H_1$, $H_2$ sedemikian
sehingga $H_1\subseteq E\subseteq H_2$ dan

\Centerline{$\mu(H_2\setminus H_1)=\mu(H_2\setminus
E)=\mu(E\setminus H_1)=0$.}

(c) Jika $A\subseteq\BbbR^r$ sembarang himpunan, maka $A$ memiliki suatu
selubung terukur yang merupakan himpunan Borel.

(d) Jika $f$ adalah fungsi bernilai real terukur Lebesgue yang didefinisikan
pada suatu subhimpunan dari $\BbbR^r$, maka terdapat himpunan Borel koterabaikan
$H\subseteq\BbbR^r$ sedemikian sehingga $f\restr H$ terukur Borel.

\proof{{\bf (a)(i)} Pertama-tama perhatikan bahwa jika
$I\subseteq\BbbR^r$ merupakan interval setengah terbuka dan $\epsilon>0$,
maka $I=\emptyset$ sudah terbuka, atau $I$ dapat dinyatakan sebagai
$\coint{a,b}$, dengan $a=(\alpha_1,\ldots,\alpha_r)$,
$b=(\beta_1,\ldots,\beta_r)$ dan $\alpha_i<\beta_i$ untuk setiap $i$.
Dalam kasus kedua, $G=\ooint{a-\epsilon(b-a),b}$ merupakan himpunan terbuka
yang memuat $I$, dan

\Centerline{$\mu G=\prod_{i=1}^r(1+\epsilon)(\beta_i-\alpha_i)
=(1+\epsilon)^r\mu I$,}

\noindent menurut rumus dalam 114G/115G.

\medskip

\quad{\bf (ii)} Sekarang, jika diberikan $\epsilon>0$, terdapat barisan
$\sequencen{I_n}$ interval setengah terbuka yang mencakup $A$ sedemikian
sehingga $\sum_{n=0}^{\infty}\mu I_n\le\mu^*A+\epsilon$. Untuk setiap
$n$, misalkan $G_n\supseteq I_n$ suatu himpunan terbuka dengan ukuran paling
besar $(1+\epsilon)^r\mu I_n$. Maka
$G=\bigcup_{n\in\Bbb N}G_n$ terbuka (1A2Bd), dan $A\subseteq G$; selain
itu,

\Centerline{$\mu G\le\sum_{n=0}^{\infty}\mu G_n
\le(1+\epsilon)^r\sum_{n=0}^{\infty}\mu I_n
\le(1+\epsilon)^r(\mu^*A+\epsilon)$.}

\noindent Karena $\epsilon$ sembarang,
$\mu^*A\ge\inf\{\mu G:G$ terbuka, $G\supseteq A\}$.

\medskip

\quad{\bf (iii)} Selanjutnya, dengan menggunakan (ii), untuk setiap
$n\in\Bbb N$ kita dapat memilih himpunan terbuka $G_n\supseteq A$
sedemikian sehingga $\mu G_n\le\mu^*A+2^{-n}$. Tetapkan
$H_0=\bigcap_{n\in\Bbb N}G_n$; maka $H_0$ merupakan himpunan Borel,
$A\subseteq H_0$, dan

\Centerline{$\mu H_0\le\inf_{n\in\Bbb N}\mu G_n\le\mu^*A$.}

\medskip

\quad{\bf (iv)} Di sisi lain, kita tentu memiliki
$\mu^*A\le\mu^*H=\mu H$ untuk setiap himpunan Borel $H\supseteq A$. Jadi
kita harus memiliki

\Centerline{$\mu^*A\le\inf\{\mu G:G\text{ terbuka},\,G\supseteq A\}$,}

\noindent dan

\Centerline{$\mu^*A=\mu H_0
=\min\{\mu H:H\text{ Borel},\,H\supseteq A\}$.}

\medskip


{\bf (b)(i)} Untuk setiap $n\in\Bbb N$, tetapkan
$E_n=E\cap B(\tbf{0},n)$. Misalkan $G_n\supseteq E_n$ suatu himpunan
terbuka dengan ukuran paling besar $\mu E_n+2^{-n}$; maka (karena
$\mu B(\tbf{0},n)<\infty$) $\mu(G_n\setminus E_n)\le 2^{-n}$.
Sekarang, untuk setiap $n$, tetapkan $G'_n=\bigcup_{m\ge n}G_m$; maka
$G'_n$ terbuka, $E=\bigcup_{m\ge n}E_m\subseteq G'_n$, dan

\Centerline{$\mu(G'_n\setminus E)\le\sum_{m=n}^{\infty}\mu(G_m\setminus
E)\le\sum_{m=n}^{\infty}\mu(G_m\setminus E_m)
\le\sum_{m=n}^{\infty}2^{-m}=2^{-n+1}$.}

\noindent Dengan menetapkan $H_2=\bigcap_{n\in\Bbb N}G'_n$, kita melihat
bahwa $H_2$ merupakan himpunan Borel yang memuat $E$ dan bahwa
$\mu(H_2\setminus E)=0$.

\medskip

\quad{\bf (ii)} Dengan mengulangi argumen (i) menggunakan
$\BbbR^r\setminus E$ sebagai pengganti $E$, kita memperoleh suatu himpunan
Borel $\tilde H_2\supseteq\BbbR^r\setminus E$ sedemikian sehingga
$\mu(\tilde H_2\setminus(\BbbR^r\setminus E))=0$; sekarang
$H_1=\BbbR^r\setminus \tilde H_2$ merupakan himpunan Borel yang termuat
dalam $E$ dan

\Centerline{$\mu(E\setminus H_1)
=\mu(\tilde H_2\setminus(\BbbR^r\setminus E))=0$.}

\noindent Tentu saja sekarang kita juga memiliki

\Centerline{$\mu(H_2\setminus H_1)=\mu(H_2\setminus E)+\mu(E\setminus
H_1)=0$.}

\medskip

\quad{\bf (iii)} Dengan kembali menggunakan gagasan (i), untuk setiap
$n\in\Bbb N$ terdapat himpunan terbuka
$\tilde G_n\supseteq B(\tbf{0},n)\setminus E$ sedemikian sehingga

\Centerline{$\mu(\tilde G_n\cap E_n)\le\mu(\tilde
G_n\setminus(B(\tbf{0},n)\setminus E))\le 2^{-n}$.}

\noindent Tetapkan

\Centerline{$F_n=B(\tbf{0},n)\setminus\tilde G_n
=B(\tbf{0},n)\cap(\BbbR^r\setminus\tilde G_n)$;}

\noindent maka $F_n$ tertutup (1A2Fd) dan berbatas, serta
$F_n\subseteq E_n\subseteq E$. Selain itu,

\Centerline{$\mu E_n=\mu F_n+\mu(E_n\setminus F_n)
=\mu F_n+\mu(\tilde G_n\cap E_n)
\le\mu F_n+2^{-n}$.}

\noindent Jadi

\Centerline{$\mu E=\lim_{n\to\infty}\mu E_n
\le\sup_{n\in\Bbb N}\mu F_n\le
\sup\{\mu F:F\text{ tertutup dan berbatas},\,F\subseteq E\}$,}

\noindent dan

\Centerline{$\mu E
=\sup\{\mu F:F\text{ tertutup dan berbatas},\,F\subseteq E\}$.}

\medskip

{\bf (c)} Misalkan $E$ sembarang selubung terukur dari $A$ (132Ef), dan
$H\supseteq E$ suatu himpunan Borel sedemikian sehingga
$\mu(H\setminus E)=0$; maka
$\mu^*(F\cap A)=\mu(F\cap E)=\mu(F\cap H)$ untuk setiap himpunan terukur
$F$, sehingga $H$ merupakan selubung terukur dari $A$.

\medskip

{\bf (d)} Tetapkan $D=\dom f$ dan tuliskan $\Cal B$ untuk
aljabar-$\sigma$ himpunan-himpunan Borel. Untuk setiap bilangan rasional
$q$, misalkan $E_q$ suatu himpunan terukur sedemikian sehingga
$\{x:f(x)\le q\}=E_q\cap D$. Misalkan $H_q$, $H_q'\in\Cal B$ sedemikian
sehingga $H_q\subseteq E_q\subseteq H'_q$ dan
$\mu(H'_q\setminus H_q)=0$. Misalkan $H$ adalah himpunan Borel
koterabaikan $\BbbR^r\setminus\bigcup(H'_q\setminus H_q)$. Maka

\Centerline{$\{x:(f\restr H)(x)\le q\}=H\cap E_q\cap D=H_q\cap D\cap H$}

\noindent termasuk dalam aljabar-$\sigma$ subruang $\Cal B(D)$ untuk
setiap $q\in\Bbb Q$. Untuk $a\in\Bbb R$ irasional, tetapkan
$H_a=\bigcap_{q\in\Bbb Q,q\ge a}H_q$; maka $H_a\in\Cal B$, dan

\Centerline{$\{x:(f\restr H)(x)\le a\}=H_a\cap\dom(f\restr H)$.}

\noindent Dengan demikian, $f\restr H$ terukur Borel.
}%end of proof of 134F

\cmmnt{\medskip

\noindent{\bf Catatan} Penekanan pada himpunan tertutup {\it berbatas}
dalam bagian (b) proposisi ini disebabkan oleh sifat-sifat topologisnya yang
penting, khususnya fakta bahwa himpunan-himpunan tersebut `kompak'. Ini
merupakan salah satu fakta terpenting tentang ukuran Lebesgue, sebagaimana
akan tampak dalam Jilid 4. Saya akan membahas `kekompakan' secara singkat
dalam \S2A2 di Jilid 2.
}%end of comment


\vleader{48pt}{134G}{Himpunan Cantor}\cmmnt{ Salah satu tujuan teori ukuran
dan integral Lebesgue adalah mempelajari himpunan dan fungsi yang lebih tak
beraturan daripada yang dapat ditangani dengan metode-metode yang lebih
primitif. Dalam beberapa paragraf berikut saya membahas himpunan dan fungsi
{\it terukur} yang, dari sudut pandang teori saat ini, dapat ditangani tanpa
menjadi sepele. Mulai sekarang, $\mu$ akan merupakan ukuran Lebesgue pada
$\Bbb R$.

\medskip

}{\bf (a)} `Himpunan Cantor' $C\subseteq[0,1]$ didefinisikan sebagai irisan
barisan himpunan $\sequencen{C_n}$, yang dibangun sebagai
berikut. $C_0=[0,1]$. Andaikan $C_n$ terdiri atas $2^n$ interval tertutup
yang saling lepas, masing-masing panjangnya $3^{-n}$; ambil setiap interval
ini dan hapus sepertiga tengahnya sehingga menghasilkan dua interval tertutup
yang masing-masing panjangnya $3^{-n-1}$; ambil $C_{n+1}$ sebagai gabungan
$2^{n+1}$ interval tertutup yang terbentuk demikian, lalu lanjutkan.
Perhatikan bahwa $\mu C_n=({2\over 3})^n$ untuk setiap $n$.

\medskip

\cmmnt{
\def\Caption{Mendekati himpunan Cantor}
\picture{mt134g}{80pt}
% created with mt134g.for
% domain -0.1  1.1  scale  0.15

\noindent}{\bf Himpunan Cantor} adalah
$C=\bigcap_{n\in\Bbb N}C_n$. Ukurannya adalah

\Centerline{$\mu C=\lim_{n\to\infty}\mu C_n
=\lim_{n\to\infty}(\Bover23)^n=0$.}

\header{134Gb}{\bf (b)} Setiap $C_n$ juga dapat digambarkan sebagai
himpunan bilangan real yang dapat dinyatakan sebagai
$\sum_{j=1}^{\infty}3^{-j}\epsilon_j$, dengan setiap $\epsilon_j$ bernilai
$0$, $1$, atau $2$, dan $\epsilon_j\ne 1$ untuk $j\le n$. Akibatnya, $C$
sendiri adalah himpunan bilangan yang dapat dinyatakan sebagai
$\sum_{j=1}^{\infty}3^{-j}\epsilon_j$, dengan setiap $\epsilon_j$ bernilai
$0$ atau $2$; yaitu, himpunan bilangan antara $0$ dan $1$ yang dapat
dinyatakan dalam bentuk terner tanpa angka 1. Representasi dalam setiap kasus
bersifat unik, sehingga kita memiliki bijeksi
$\phi:\{0,1\}^{\Bbb N}\to C$ yang didefinisikan dengan menuliskan

\Centerline{$\phi(z)=\Bover23\sum_{j=0}^{\infty}3^{-j}z(j)$}

\noindent untuk setiap $z\in\{0,1\}^{\Bbb N}$.

\leader{134H}{Fungsi Cantor}\cmmnt{ Melanjutkan dari 134G, kita memperoleh
konstruksi berikut.


\medskip

}{\bf (a)} Untuk setiap $n\in\Bbb N$\cmmnt{ kita} mendefinisikan fungsi
$f_n:[0,1]\to[0,1]$ dengan menetapkan

\Centerline{$f_n(x)=(\Bover32)^n\mu(C_n\cap[0,x])$}

\noindent untuk setiap $x\in [0,1]$. \cmmnt{Karena $C_n$ hanyalah suatu
gabungan berhingga interval,} $f_n$ merupakan fungsi poligonal, dengan
$f_n(0)=0$, $f_n(1)=1$; $f_n$ konstan pada masing-masing dari $2^n-1$ interval
terbuka penyusun $[0,1]\setminus C_n$, dan naik dengan kemiringan
$({3\over 2})^n$ pada masing-masing dari $2^n$ interval tertutup penyusun $C_n$.

\cmmnt{
\def\Caption{Mendekati fungsi Cantor:
fungsi-fungsi $f_0$, $f_1$, $f_2$, $\pmb{f_3}$}
\picture{mt134ha1}{195pt}
% computed with mt134ha1.for
% domain -0.1  1.1  scale  0.15
%  dash/gap   0.3/0.3  0.267/0.233  0.233/0.167  0.2/0.1
}%end of comment

\cmmnt{Jika interval ke-$j$ dari $C_n$, dihitung dari kiri, adalah
$[a_{nj},b_{nj}]$, maka $f_n(a_{nj})=2^{-n}(j-1)$ dan
$f_n(b_{nj})=2^{-n}j$. Selain itu, $a_{nj}=a_{n+1,2j-1}$ dan
$b_{nj}=b_{n+1,2j}$; dengan demikian, atau melalui cara lain,
$f_{n+1}(a_{nj})=f_n(a_{nj})$ dan
$f_{n+1}(b_{nj})=f_{n}(b_{nj})$, dan $f_{n+1}$ berimpit dengan $f_n$
pada semua titik ujung interval-interval $C_n$, dan karenanya pada
$[0,1]\setminus C_n$.

Di dalam setiap interval tertentu $[a_{nj},b_{nj}]$ dari $C_n$, selisih
terbesar antara $f_n(x)$ dan $f_{n+1}(x)$ terjadi pada titik-titik ujung
baru di dalam interval tersebut, yaitu $b_{n+1,2j-1}$ dan $a_{n+1,2j}$;
dan besar selisihnya adalah ${1\over 6}2^{-n}$ (sebab, misalnya,
$f_n(b_{n+1,2j-1})={2\over 3}f_n(a_{nj})+{1\over 3}f_n(b_{nj})$, sedangkan
$f_{n+1}(b_{n+1,2j-1})={1\over 2}f_n(a_{nj})+{1\over 2}f_n(b_{nj})$).
Dengan demikian kita memiliki
$|f_{n+1}(x)-f_n(x)|\le{1\over 6}2^{-n}$ untuk setiap $n\in\Bbb N$,
$x\in[0,1]$. Karena $\sum_{n=0}^{\infty}{1\over 6}2^{-n}<\infty$,
}%end of comment
$\sequencen{f_n}$ konvergen seragam menuju suatu fungsi
$f:[0,1]\to[0,1]$, dan $f$ kontinu. $f$ adalah {\bf fungsi Cantor} atau
{\bf Tangga Setan}.

\def\Caption{Fungsi Cantor}
\picture{mt134ha2}{187.43pt}
% computed with  mt134ha2.for
% domain  -0.1  1.1  scale  0.15



\header{134Hb}{\bf (b)} Karena setiap $f_n$ tak-menurun, $f$ juga
demikian. \cmmnt{Jika $x\in[0,1]\setminus C$, terdapat suatu $n$ sedemikian
sehingga $x\in[0,1]\setminus C_n$; misalkan $I$ interval terbuka dari
$[0,1]\setminus C_n$ yang memuat $x$; maka $f_{m+1}$ berimpit pada $I$
dengan $f_m$ untuk setiap $m\ge n$, sehingga $f$ berimpit pada $I$ dengan
$f_n$, dan $f$ konstan pada $I$. Dengan demikian, khususnya, turunan $f'(x)$
ada dan bernilai $0$ untuk setiap $x\in[0,1]\setminus C$; sehingga}
$f'$ bernilai nol hampir di mana-mana dalam $[0,1]$. \cmmnt{Tentu saja,
$f(0)=0$ dan $f(1)=1$, karena $f_n(0)=0$, $f_n(1)=1$ untuk setiap $n$.
Akibatnya,} $f:[0,1]\to[0,1]$ surjektif\cmmnt{ (menurut Teorema Nilai
Antara)}.

\header{134Hc}{\bf (c)} Misalkan $\phi:\{0,1\}^{\Bbb N}\to C$ adalah
fungsi yang dijelaskan dalam 134Gb. Maka
$f(\phi(z))=\bover12\sum_{j=0}^{\infty}2^{-j}z(j)$ untuk setiap
$z\in\{0,1\}^{\Bbb N}$. \prooflet{\Prf\ Tetapkan
$z=(\zeta_0,\zeta_1,\zeta_2,\ldots)$ dalam $\{0,1\}^{\Bbb N}$, dan untuk
setiap $n$ ambil $I_n$ sebagai interval komponen $C_n$ yang memuat
$\phi(z)$. Maka $I_{n+1}$ merupakan sepertiga kiri $I_n$ jika
$\zeta_n=0$, dan sepertiga kanannya jika $\zeta_n=1$. Dengan mengambil
$a_n$ sebagai titik ujung kiri $I_n$, kita melihat bahwa

\Centerline{$a_{n+1}=a_n+\Bover233^{-n}\zeta_n$,
\quad$f_{n+1}(a_{n+1})=f_n(a_n)+\Bover122^{-n}\zeta_n$}

\noindent untuk setiap $n$. Sekarang

\Centerline{$\phi(z)=\lim_{n\to\infty}a_n$,
\quad$f(\phi(z))=\lim_{n\to\infty}f(a_n)=\lim_{n\to\infty}f_n(a_n)
=\Bover12\sum_{j=0}^{\infty}2^{-j}\zeta_j$,}

\noindent seperti yang dinyatakan.\ \Qed}%end of prooflet
\cmmnt{

Khususnya,} $f[C]=[0,1]$. \prooflet{\Prf\ Setiap $x\in[0,1]$ dapat
dinyatakan sebagai $\sum_{j=0}^{\infty}2^{-j-1}z(j)=f(\phi(z))$ untuk
suatu $z\in\{0,1\}^{\Bbb N}$.\ \Qed}

\leader{134I}{Fungsi Cantor yang dimodifikasi}\cmmnt{ Saya melanjutkan
argumen 134G-134H.

\medskip

}{\bf (a)} Tinjau rumus

\Centerline{$g(x)=\Bover12(x+f(x))$,}

\noindent dengan $f$ fungsi Cantor, sebagaimana didefinisikan dalam 134H;
rumus ini mendefinisikan suatu fungsi kontinu $g:[0,1]\to[0,1]$ yang naik
ketat\cmmnt{ (karena $f$ tak-menurun)} serta memenuhi $g(0)=0$, $g(1)=1$;
\cmmnt{akibatnya, menurut Teorema Nilai Antara,} $g$ bijektif, dan inversnya
$g^{-1}:[0,1]\to[0,1]$ kontinu.
\cmmnt{

Sekarang} $g[C]$ merupakan himpunan tertutup dan
$\mu g[C]=\bover12$. \prooflet{\Prf\ Karena $g$ merupakan permutasi
titik-titik $[0,1]$,
$[0,1]\setminus g[C]=g[\,[0,1]\setminus C]$.
Untuk setiap interval terbuka $I_{nj}=\ooint{b_{nj},a_{n,j+1}}$ yang
menyusun $[0,1]\setminus C_n$, kita melihat bahwa
$g[I_{nj}]=\ooint{g(b_{nj}),g(a_{n,j+1})}$ memiliki panjang tepat separuh
panjang $I_{nj}$. Akibatnya
$g[\,[0,1]\setminus C]=\bigcup_{n\ge
1,1\le j<2^n}g[I_{nj}]$ terbuka, dan

$$\eqalign{\mu(g[\,[0,1]\setminus C_n])
&=\sum_{j=1}^{2^n-1}g(a_{n,j+1})-g(b_{nj})
=\Bover12\sum_{j=1}^{2^n-1}a_{n,j+1}-b_{nj}\cr
&=\Bover12\mu([0,1]\setminus C_n)
=\Bover12(1-(\Bover23)^n)\cr}$$

\noindent (134Ga). Karena $\sequencen{[0,1]\setminus C_n}$ merupakan
barisan himpunan menaik dengan gabungan $[0,1]\setminus C$,

\Centerline{$\mu g([\,[0,1]\setminus C])
=\lim_{n\to\infty}\mu g([\,[0,1]\setminus C_n])
=\Bover12$.}

\noindent Jadi $g[C]=[0,1]\setminus g[\,[0,1]\setminus C]$ tertutup dan
$\mu g[C]=\bover12$. \Qed}

\header{134Ib}{\bf (b)} Menurut 134D terdapat suatu himpunan
$D\subseteq\Bbb R$ sedemikian sehingga

\Centerline{$\mu^*(g[C]\cap D)=\mu^*(g[C]\setminus D)
=\mu g[C]=\Bover12$;}

\noindent tetapkan $A=g[C]\cap D$. \cmmnt{Tentu saja} $A$ tidak mungkin
terukur\cmmnt{, karena $\mu^*A+\mu^*(g[C]\setminus A)>\mu g[C]$}.
Namun, $g^{-1}[A]\subseteq C$ harus terukur\cmmnt{, karena $\mu^*C=0$}.
Ini berarti bahwa jika kita menetapkan
$h=\chi(g^{-1}[A]):[0,1]\to\Bbb R$, maka $h$ terukur; tetapi
$hg^{-1}\cmmnt{\mskip5mu =\chi A:[0,1]\to\Bbb R}$ tidak terukur.

Jadi, {\bf komposisi suatu fungsi terukur dengan fungsi kontinu tidak harus
terukur}. \cmmnt{Bandingkan ini dengan 121Eg.}

\leader{134J}{Contoh-contoh lain}\cmmnt{ Menurut saya, ada baiknya
menyediakan ruang untuk menguraikan secara terperinci dua contoh dasar lain
mengenai himpunan-himpunan terukur Lebesgue.

\medskip

}{\bf (a)}\cmmnt{ Sebagaimana telah diamati dalam 114G, setiap subhimpunan
terhitung dari $\Bbb R$ terabaikan. Khususnya, $\Bbb Q$ terabaikan (111Eb).
Kita dapat mengatakan lebih jauh.} Misalkan $\sequencen{q_n}$ suatu barisan
yang menjangkau $\Bbb Q$, dan untuk setiap $n\in\Bbb N$ tetapkan

\Centerline{$I_n=\ooint{q_n-2^{-n},q_n+2^{-n}}$,}

\Centerline{$G_n=\bigcup_{k\ge n}I_k$.}

\noindent Maka $G_n$ merupakan himpunan terbuka berukuran paling besar
$\sum_{k=n}^{\infty}2\cdot 2^{-k}=4\cdot 2^{-n}$, dan memuat semua titik
$\Bbb Q$ kecuali sejumlah berhingga, sehingga rapat\cmmnt{ (yaitu, beririsan
dengan setiap interval tak-sepele)}. Tetapkan $F_n=\Bbb R\setminus G_n$;
maka $F_n$ tertutup, $\mu(\Bbb R\setminus F_n)\le 4/2^n$, tetapi $F_n$
tidak\cmmnt{ memuat satu pun $q_k$ dengan $k\ge n$, sehingga $F_n$ tidak
mungkin} memuat interval tak-sepele apa pun.
\cmmnt{Perhatikan bahwa $\sequencen{G_n}$ tak-menaik, sehingga}
$\sequencen{F_n}$ tak-menurun.

\header{134Jb}{\bf (b)}\cmmnt{ Kita dapat mengembangkan konstruksi di atas
sebagai berikut.} Terdapat suatu himpunan terukur $E\subseteq\Bbb R$
sedemikian sehingga $\mu(I\cap E)>0$ dan $\mu(I\setminus E)>0$ untuk setiap
interval tak-sepele $I\subseteq\Bbb R$. \prooflet{\Prf\ Pertama-tama
perhatikan bahwa jika $k$, $n\in\Bbb N$, terdapat $j\ge n$ sedemikian
sehingga $q_j\in I_k$, sehingga $I_k\cap I_j\ne\emptyset$ dan
$\mu(I_k\setminus F_n)>0$. Sekarang pasti terdapat suatu $l>n$ sedemikian
sehingga $\mu G_l<\mu(I_k\setminus F_n)$, sehingga

\Centerline{$\mu(I_k\cap F_l\setminus F_n)
=\mu((I_k\setminus F_n)\setminus G_l)>0$.}

Pilih $n_0<n_1<n_2<\ldots$ sebagai berikut. Mulailah dengan $n_0=0$.
Jika $n_{2k}$ telah diberikan, dengan $k\in\Bbb N$, pilih $n_{2k+1}$,
$n_{2k+2}$ sedemikian sehingga

\Centerline{$\mu(I_k\cap F_{n_{2k+1}}\setminus F_{n_{2k}})>0,\quad
\mu(I_k\cap F_{n_{2k+2}}\setminus F_{n_{2k+1}})>0$.}

\noindent Lanjutkan.

Setelah induksi selesai, tetapkan

\Centerline{$E=\bigcup_{k\in\Bbb N}F_{n_{2k+1}}\setminus F_{n_{2k}}$,
\quad$H=\bigcup_{k\in\Bbb N}F_{n_{2k+2}}\setminus F_{n_{2k+1}}$.}

\noindent Karena $\sequence{k}{F_k}$ tak-menurun,
$E\cap H=\emptyset$. Jika $k\in\Bbb N$, $E\cap I_k$ dan $H\cap I_k$
keduanya mempunyai ukuran positif.

Sekarang andaikan $I\subseteq\Bbb R$ suatu interval dengan lebih dari satu
titik; andaikan $a$, $b\in I$ dan $a<b$. Maka terdapat $m\in\Bbb N$
sedemikian sehingga $4\cdot 2^{-m}\le b-a$; sekarang terdapat suatu
$k\ge m$ sedemikian sehingga $q_k\in[a+2^{-m},b-2^{-m}]$, sehingga
$I_k\subseteq I$ dan

\Centerline{$\mu(I\cap E)\ge\mu(E\cap I_k)>0$,
\quad$\mu(I\setminus E)\ge\mu(H\cap I_k)>0$.  \Qed}}

\header{134Jc}{\bf (c)}
\cmmnt{Ini menunjukkan bahwa} $E$ dan komplemennya adalah himpunan-himpunan
terukur yang\cmmnt{ bukan sekadar sama-sama rapat (seperti $\Bbb Q$ dan
$\Bbb R\setminus\Bbb Q$), melainkan} rapat `secara esensial' dalam arti
bahwa keduanya beririsan dengan setiap interval terbuka tak-kosong dalam
suatu himpunan berukuran positif, sehingga\cmmnt{ (misalnya)} $E\setminus A$
rapat untuk setiap himpunan terabaikan $A$.

\leader{*134K}{Integral Riemann}\cmmnt{ Dalam menulis buku ini, saya telah
berusaha mengandaikan sesedikit mungkin pengetahuan sebelumnya. Khususnya,
Anda tidak perlu telah mempelajari integral Riemann. Namun demikian, jika
Anda telah mempelajari teori dasar integral Riemann --- yang sesungguhnya
bukan hanya latihan sangat baik dalam teknik analisis $\epsilon$-$\delta$,
melainkan juga sumber gagasan yang terus-menerus bagi bidang ini --- saya
harap Anda ingin menghubungkannya dengan materi yang sedang kita pelajari;
baik karena Anda tidak ingin merasa bahwa jerih payah Anda sia-sia, maupun
karena Anda barangkali telah mengembangkan sejumlah intuisi yang akan tetap
berharga apabila disesuaikan dengan tepat pada konteks baru. Oleh karena itu,
saya memberikan uraian singkat mengenai hubungan antara metode integral
Riemann dan Lebesgue pada garis bilangan real.

\medskip

{\bf (a)} Ada banyak cara untuk menggambarkan integral Riemann; saya memilih
salah satu yang populer. Jika $[a,b]$ adalah interval tertutup tak-sepele
dalam $\Bbb R$, saya mengatakan bahwa suatu {\bf pembagian} atas $[a,b]$
adalah daftar berhingga $D=(a_0,a_1,\ldots,a_n)$, dengan $n\ge 1$,
sedemikian sehingga $a=a_0<a_1<\ldots< a_n=b$. Jika sekarang $f$ merupakan
fungsi bernilai real yang didefinisikan (setidaknya) pada $[a,b]$ dan
berbatas pada $[a,b]$, {\bf jumlah atas} dan {\bf jumlah bawah} dari $f$
pada $[a,b]$ yang diturunkan dari $D$ adalah

\Centerline{$S_D(f)
=\sum_{i=1}^n(a_i-a_{i-1})\sup_{x\in\ooint{a_{i-1},a_i}}f(x)$,}

\Centerline{$s_D(f)
=\sum_{i=1}^n(a_i-a_{i-1})\inf_{x\in\ooint{a_{i-1},a_i}}f(x)$.}

\noindent Anda perlu membuktikan bahwa jika $D$ dan $D'$ merupakan dua
pembagian atas $[a,b]$, maka $s_D(f)\le S_{D'}(f)$. Sekarang definisikan
{\bf integral Riemann atas} dan {\bf integral Riemann bawah} dari $f$
sebagai

\Centerline{$U_{[a,b]}(f)=\inf\{S_D(f):D$ pembagian atas $[a,b]\}$,}

\Centerline{$L_{[a,b]}(f)=\sup\{s_D(f):D$ pembagian atas $[a,b]\}$.}

\noindent Periksa bahwa $L_{[a,b]}(f)$ pasti lebih kecil daripada atau sama
dengan $U_{[a,b]}(f)$. Akhirnya, nyatakan $f$ {\bf terintegralkan Riemann
pada} $[a,b]$ jika $U_{[a,b]}(f)=L_{[a,b]}(f)$, dan dalam kasus ini ambil
nilai bersama tersebut sebagai {\bf integral Riemann} $\Rint_a^bf$ dari
$f$ pada $[a,b]$.

\spheader 134Kb} %end of comment
Jika $f:[a,b]\to\Bbb R$ terintegralkan Riemann, maka fungsi tersebut terintegralkan
Lebesgue, dengan integral yang sama. \prooflet{\Prf\ Untuk setiap pembagian
$D=(a_0,\ldots,a_n)$ atas $[a,b]$, definisikan
$g_D$, $h_D:[a,b]\to\Bbb R$ dengan menetapkan

\Centerline{$g_D(x)=\inf\{f(y):y\in\ooint{a_{i-1},a_i}\}$ jika
$a_{i-1}<x<a_i$, \quad $g_D(a_i)=f(a_i)$ untuk setiap $i$,}

\Centerline{$h_D(x)=\sup\{f(y):y\in\ooint{a_{i-1},a_i}\}$ jika
$a_{i-1}<x<a_i$, \quad $h_D(a_i)=f(a_i)$ untuk setiap $i$.}

\noindent Maka $g_D$ dan $h_D$ konstan pada setiap interval
$\ooint{a_{i-1},a_i}$, sehingga semua himpunan $\{x:g_D(x)\le c\}$,
$\{x:h_D(x)\le c\}$ merupakan gabungan berhingga interval-interval, dan
$g_D$ serta $h_D$ terukur; selain itu,

\Centerline{$\int g_Dd\mu
=s_D(f)$,\quad $\int h_Dd\mu=S_D(f)$.}

\noindent Akibatnya,

$$\eqalign{\Rint_a^bf=L_{[a,b]}(f)&=\sup_D\int g_Dd\mu
\le\underline{\int}fd\mu\cr
&\le\overline{\int}fd\mu
\le\inf_D\int h_Dd\mu=U_{[a,b]}(f)=\Rint_a^bf,\cr}$$

\noindent dan $\overline{\int}fd\mu=\underline{\int}fd\mu=\Rint_a^bf$,
sehingga $\int fd\mu$ ada dan sama dengan $\Rint_a^bf$ (133Jd).\ \Qed
}%end of prooflet


\cmmnt{\header{134Kc}{\bf (c)} Pembahasan di atas menyangkut integral
Riemann `wajar', untuk fungsi berbatas pada interval berbatas. Untuk fungsi
tak-berbatas dan interval tak-berbatas, digunakan berbagai bentuk integral
`tak wajar'; misalnya, integral Riemann tak wajar
$\int_0^{\infty}\bover{\sin x}{x}dx$ diambil sebagai
$\lim_{a\to\infty}\int_0^a\bover{\sin x}{x}dx$, sedangkan
$\int_0^1\ln x\,dx$ diambil sebagai
$\lim_{a\downarrow 0}\int_a^1\ln x\,dx$.
Di antara keduanya, integral kedua ada sebagai integral Lebesgue, tetapi yang
pertama tidak, karena $\int_0^{\infty}|\bover{\sin x}{x}|dx=\infty$.
Kemampuan integral Lebesgue untuk menangani secara langsung fungsi
tak-berbatas `terintegralkan secara mutlak' pada domain tak-berbatas berarti
bahwa fungsi yang dapat disebut `terintegralkan secara bersyarat' terdorong
ke latar belakang teori. Dalam Bab 48 Jilid 4 saya akan membahas teori umum
fungsi-fungsi semacam itu, tetapi untuk sementara waktu saya akan
menanganinya satu per satu pada kesempatan langka ketika fungsi tersebut
muncul.
}%end of comment

\leader{*134L}{}\cmmnt{ Sebenarnya terdapat suatu karakterisasi indah bagi
fungsi-fungsi terintegralkan Riemann, sebagai berikut.

\medskip

\noindent}{\bf Proposisi} Jika $a<b$ dalam $\Bbb R$, suatu fungsi berbatas
$f:[a,b]\to\Bbb R$ terintegralkan Riemann jika dan hanya jika kontinu
hampir di mana-mana dalam $[a,b]$.

\proof{{\bf (a)} Andaikan $f$ terintegralkan Riemann. Untuk setiap
$x\in[a,b]$, tetapkan

\Centerline{$g(x)
=\sup_{\delta>0}\inf_{y\in[a,b],|y-x|\le\delta}f(y)$,}

\Centerline{$h(x)
=\inf_{\delta>0}\sup_{y\in[a,b],|y-x|\le\delta}f(y)$,}

\noindent sehingga $f$ kontinu di $x$ jika dan hanya jika $g(x)=h(x)$.
Kita memiliki $g\le f\le h$, sehingga jika $D$ sembarang pembagian atas
$[a,b]$, maka $S_D(g)\le S_D(f)\le S_D(h)$ dan
$s_D(g)\le s_D(f)\le s_D(h)$. Namun sebenarnya $S_D(f)=S_D(h)$ dan
$s_D(g)=s_D(f)$, karena pada setiap interval terbuka
$\ooint{c,d}\subseteq[a,b]$ kita harus memiliki

\Centerline{$\inf_{x\in\ooint{c,d}}g(x)=\inf_{x\in\ooint{c,d}}f(x)$,
\quad
$\sup_{x\in\ooint{c,d}}f(x)=\sup_{x\in\ooint{c,d}}h(x)$.}

\noindent Akibatnya,

\Centerline{$L_{[a,b]}(f)=L_{[a,b]}(g)\le U_{[a,b]}(g)\le
U_{[a,b]}(f)$,}

\Centerline{$L_{[a,b]}(f)\le L_{[a,b]}(h)\le U_{[a,b]}(h)=
U_{[a,b]}(f)$.}

Karena $f$ terintegralkan Riemann, $g$ dan $h$ keduanya harus terintegralkan
Riemann, dengan integral sama dengan $\Rint_a^bf$. Menurut 134Kb, keduanya
terintegralkan Lebesgue, dengan integral yang sama. Namun $g\le h$, sehingga
$g\eae h$, menurut 122Rd. Sekarang $f$ kontinu di setiap titik tempat $g$
dan $h$ berimpit, sehingga $f$ kontinu hampir di mana-mana.

\medskip

{\bf (b)} Sekarang andaikan $f$ kontinu hampir di mana-mana. Untuk setiap $n\in\Bbb N$,
misalkan $D_n$ pembagian atas $[a,b]$ menjadi $2^n$ bagian sama panjang.
Tetapkan

\Centerline{$h_n(x)=\sup_{y\in\ooint{c,d}}f(y)$,
\quad$g_n(x)=\inf_{y\in\ooint{c,d}}f(y)$}

\noindent jika $\ooint{c,d}$ merupakan interval terbuka dari $D_n$ yang
memuat $x$; sebagai konvensi, tetapkan $h_n(x)=g_n(x)=f(x)$ jika $x$
merupakan salah satu titik dalam daftar $D_n$. Maka $\sequencen{g_n}$
merupakan barisan fungsi menaik, sedangkan $\sequencen{h_n}$ merupakan
barisan fungsi menurun; setiap fungsinya konstan pada masing-masing anggota suatu keluarga
berhingga interval yang mencakup $[a,b]$; dan
$s_{D_n}(f)=\int g_nd\mu$, $S_{D_n}(f)=\int h_nd\mu$. Selanjutnya,

\Centerline{$\lim_{n\to\infty}g_n(x)=\lim_{n\to\infty}h_n(x)=f(x)$}

\noindent pada setiap titik $x$ tempat $f$ kontinu; sehingga
$f\eae\lim_{n\to\infty}g_n\eae\lim_{n\to\infty}h_n$. Menurut Teorema
Konvergensi Terdominasi Lebesgue (123C),

\Centerline{$\lim_{n\to\infty}\int g_nd\mu
=\int fd\mu=\lim_{n\to\infty}\int h_nd\mu$;}

\noindent tetapi ini berarti bahwa

\Centerline{$L_{[a,b]}(f)\ge\int fd\mu\ge U_{[a,b]}(f)$,}

\noindent sehingga semua ini sama dan $f$ terintegralkan Riemann.
}%end of proof of 134L


\exercises{
\leader{134X}{Latihan dasar $\pmb{>}$(a)}
%\spheader 134Xa
Tunjukkan bahwa jika $f$ merupakan fungsi bernilai real terintegralkan pada
$\BbbR^r$, maka $\int f(x+a)dx$ ada dan sama dengan $\int f$ untuk setiap
$a\in\BbbR^r$. \Hint{mulailah dengan fungsi sederhana $f$.}
%134A

\spheader 134Xb Secara lebih umum, tunjukkan bahwa jika
$E\subseteq\BbbR^r$ terukur dan $f$ merupakan fungsi bernilai real yang
terintegralkan pada $E$ dalam arti 131D, maka
$\int_{E-a}f(x+a)dx$ ada dan sama dengan $\int_Ef$ untuk setiap
$a\in\BbbR^r$.
%134Xa, 134A

\spheader 134Xc Tunjukkan bahwa jika $C\subseteq\Bbb R$ sembarang himpunan
tak-terabaikan, maka himpunan tersebut memiliki suatu subhimpunan tak-terukur.
\Hint{gunakan metode 134B, dengan mengambil relasi $\sim$ pada subhimpunan
berbatas yang sesuai dari $C$ sebagai pengganti
$\coint{\tbf{0},\tbf{1}}$.}
%134B

\sqheader 134Xd Misalkan $\nu_g$ suatu ukuran Lebesgue–Stieltjes pada
$\Bbb R$, yang dibangun seperti dalam 114Xa dari fungsi tak-menurun
$g:\Bbb R\to\Bbb R$, dan $\Sigma_g$ domainnya. (Lihat juga 132Xg.)
Tunjukkan bahwa

\quad(i) jika $A\subseteq\Bbb R$ sembarang himpunan, maka

$$\eqalign{\nu_g^*A&=\inf\{\nu_g G:G\text{ terbuka},\,G\supseteq A\}\cr
&=\min\{\nu_g H:H\text{ Borel},\,H\supseteq A\};\cr}$$

\quad(ii) jika $E\in\Sigma_g$, maka

\Centerline{$\nu_g E=\sup\{\nu_g F:F
\text{ tertutup dan berbatas},\,F\subseteq E\}$,}

\noindent dan terdapat himpunan-himpunan Borel $H_1$, $H_2$ sedemikian
sehingga $H_1\subseteq E\subseteq H_2$ dan
$\nu_g(H_2\setminus H_1)=\nu_g(H_2\setminus E)=\nu_g(E\setminus H_1)=0$;

\quad(iii) jika $A\subseteq\Bbb R$ sembarang himpunan, maka $A$ memiliki
suatu selubung terukur yang merupakan himpunan Borel;

\quad(iv) jika $f$ adalah fungsi bernilai real $\Sigma_g$-terukur yang
didefinisikan pada suatu subhimpunan dari $\Bbb R$, maka terdapat himpunan Borel
$\nu_g$-koterabaikan $H\subseteq\Bbb R$ sedemikian sehingga $f\restr H$
terukur Borel.
%134F

\spheader 134Xe Misalkan $E\subseteq\BbbR^r$ suatu himpunan terukur, dan
$\epsilon>0$. (i) Tunjukkan bahwa terdapat suatu himpunan terbuka
$G\supseteq E$ sedemikian sehingga $\mu(G\setminus E)\le\epsilon$.
\Hint{terapkan 134Fa pada setiap himpunan $E\cap B(\tbf{0},n)$.}
(ii) Tunjukkan bahwa terdapat suatu himpunan tertutup $F\subseteq E$
sedemikian sehingga $\mu(E\setminus F)\le\epsilon$.
%134F

\spheader 134Xf Misalkan $C\subseteq[0,1]$ himpunan Cantor. Tunjukkan bahwa
$\{x+y:x$, $y\in C\}=[0,2]$ dan $\{x-y:x$, $y\in C\}=[-1,1]$.
%134G

\spheader 134Xg Misalkan $f$, $g$ fungsi-fungsi dari $\Bbb R$ ke dirinya
sendiri. Tunjukkan bahwa (i) jika $f$ dan $g$ keduanya terukur Borel, maka
komposisinya $fg$ juga demikian (ii) jika $f$ terukur Borel dan $g$ terukur
Lebesgue, maka $fg$ terukur Lebesgue (iii) jika $f$ terukur Lebesgue dan $g$
terukur Borel, maka $fg$ tidak harus terukur Lebesgue.
%134I

\spheader 134Xh Tunjukkan bahwa untuk setiap bilangan bulat $r\ge 1$
terdapat suatu himpunan terukur $E\subseteq\BbbR^r$ sedemikian sehingga
$E$ dan $\BbbR^r\setminus E$ keduanya beririsan dengan setiap interval
terbuka tak-kosong dalam suatu himpunan berukuran positif.
%134J

\spheader 134Xi Bekali $[0,1]$ dengan ukuran subruangnya. (i) Tunjukkan bahwa
terdapat suatu barisan saling lepas $\sequencen{A_n}$ yang terdiri atas
subhimpunan-subhimpunan dari $[0,1]$ yang semuanya berukuran luar $1$.
(ii) Tunjukkan bahwa terdapat suatu fungsi $f:[0,1]\to\ooint{0,1}$
sedemikian sehingga $\underline{\int}f=0$ dan $\overline{\int}f=1$.
%134B

\spheader 134Xj Misalkan $f$ suatu fungsi real terukur dan $g$ suatu fungsi
real sedemikian sehingga $\dom g\setminus\dom f$ dan
$\{x:x\in\dom g\cap\dom f$, $g(x)\ne f(x)\}$ keduanya terabaikan.
Tunjukkan bahwa $g$ terukur.
%134Xg 134I

\leader{134Y}{Latihan lanjutan (a)}
%\spheader 134Ya
Tetapkan $c>0$. Untuk $A\subseteq\BbbR^r$ tetapkan
$cA=\{cx:x\in A\}$. (i) Tunjukkan bahwa $\mu^*(cA)=c^r\mu^*A$ untuk setiap
$A\subseteq\BbbR^r$. (ii) Tunjukkan bahwa $cE$ terukur untuk setiap
$E\subseteq\BbbR^r$ terukur.
%134A
%this is in 115Xe;  can be dropped here query

\spheader 134Yb Misalkan $\langle f_{mn}\rangle_{m,n\in\Bbb N}$,
$\sequence{m}{f_m}$, $f$ merupakan fungsi-fungsi terukur bernilai real yang
didefinisikan hampir di mana-mana dalam $\BbbR^r$ dan sedemikian sehingga
$f_m\eae\lim_{n\to\infty}f_{mn}$ untuk setiap $m$ dan
$f\eae\lim_{m\to\infty}f_m$. Tunjukkan bahwa terdapat suatu barisan
$\sequence{k}{n_k}$ sedemikian sehingga
$f\eae\lim_{k\to\infty}f_{k,n_k}$.
\Hint{ambil $n_k$ sedemikian sehingga ukuran
$\{x:\|x\|\le k,\,|f_k(x)-f_{k,n_k}(x)|\ge 2^{-k}\}$ paling besar $2^{-k}$
untuk setiap $k$.}

\spheader 134Yc Misalkan $f$ suatu fungsi bernilai real terukur yang
didefinisikan hampir di mana-mana dalam $\BbbR^r$. Tunjukkan bahwa terdapat
suatu barisan $\sequencen{f_n}$ fungsi-fungsi kontinu yang konvergen menuju
$f$ hampir di mana-mana. \Hint{Tangani secara berturut-turut kasus-kasus
(i) $f=\chi I$ dengan $I$ suatu interval setengah terbuka
(ii) $f=\chi(\bigcup_{j\le n}I_j)$ dengan $I_0,\ldots,I_n$ interval-interval
setengah terbuka yang saling lepas (iii) $f=\chi E$ dengan $E$ suatu
himpunan terukur berukuran berhingga (iv) $f$ suatu fungsi sederhana
(v) $f$ umum, dengan menggunakan 134Yb pada langkah (iii) dan (v).}
%134Yb, 134F

\spheader 134Yd Misalkan $f$ suatu fungsi bernilai real yang didefinisikan
pada suatu subhimpunan dari $\BbbR^r$. Tunjukkan bahwa pernyataan-pernyataan
berikut ekuivalen: (i) $f$ terukur (ii) jika $E\subseteq\BbbR^r$ terukur
dan $\mu E>0$, terdapat himpunan terukur $F\subseteq E$ sedemikian sehingga
$\mu F>0$ dan $f\restr F$ kontinu (iii) jika $E\subseteq\BbbR^r$ terukur
dan $\gamma<\mu E$, terdapat suatu $F\subseteq E$ terukur sedemikian
sehingga $\mu F\ge\gamma$ dan $f\restr F$ kontinu. \Hint{untuk
(i)$\Rightarrow$(iii), gunakan 134Yc dan 131Ya; untuk (ii)$\Rightarrow$(i)
gunakan 121D.
Ini merupakan suatu versi {\bf teorema Lusin}.}
%134F, 134Yc

\spheader 134Ye Misalkan $\nu$ suatu ukuran pada $\Bbb R$ yang invarian
terhadap translasi dalam arti 134Ab, dan sedemikian sehingga $\nu[0,1]$
terdefinisi dan sama dengan $1$. Tunjukkan bahwa $\nu$ berimpit dengan
ukuran Lebesgue pada himpunan-himpunan Borel dalam $\Bbb R$. \Hint{Tunjukkan
terlebih dahulu bahwa $[a,1]$ termasuk dalam domain $\nu$ untuk setiap
$a\in[0,1]$, dan karenanya bahwa setiap interval setengah terbuka dengan
panjang paling besar $1$ termasuk dalam domain $\nu$; tunjukkan bahwa
$\nu\coint{a,a+2^{-n}}=2^{-n}$ untuk setiap $a\in\Bbb R$,
$n\in\Bbb N$, dan karenanya bahwa $\nu\coint{a,b}=b-a$ setiap kali $a<b$.}
%134A, 134F

\spheader 134Yf Misalkan $\nu$ suatu ukuran pada $\BbbR^r$ yang invarian
terhadap translasi dalam arti 134Ab, dengan $r>1$, dan sedemikian sehingga
$\nu[\tbf{0},\tbf{1}]$ terdefinisi dan sama dengan $1$. Tunjukkan bahwa
$\nu$ berimpit dengan ukuran Lebesgue pada himpunan-himpunan Borel dalam
$\BbbR^r$.
%134Ye, 134F

\spheader 134Yg Tunjukkan bahwa jika $f$ sembarang fungsi bernilai real yang
terintegralkan pada $\Bbb R$, dan $\epsilon>0$, terdapat suatu fungsi
kontinu $g:\Bbb R\to\Bbb R$ sedemikian sehingga $\{x:g(x)\ne 0\}$ berbatas
dan $\int|f-g|\le\epsilon$. \Hint{tunjukkan bahwa himpunan $\Phi$ dari
fungsi-fungsi $f$ dengan sifat ini memenuhi syarat-syarat 122Yb.}
%134F

\spheader 134Yh Ulangi 134Yg untuk fungsi-fungsi bernilai real yang
terintegralkan pada $\BbbR^r$, dengan $r>1$.
%134Yg, 134F

\spheader 134Yi Ulangi 134Fd, 134Xa, 134Xb, 134Yb, 134Yc, 134Yd, 134Yg
dan 134Yh untuk fungsi-fungsi bernilai kompleks.
%134Yh, 134F


\spheader 134Yj Tunjukkan bahwa jika $G\subseteq\BbbR^r$ terbuka dan
tak-kosong, maka himpunan tersebut dapat dinyatakan sebagai gabungan saling lepas suatu
barisan interval setengah terbuka yang masing-masing berbentuk
$\{x:2^{-m}n_i\le\xi_i<2^{-m}(n_i+1)$ untuk setiap $i\le r\}$, dengan
$m\in\Bbb N$, $n_1,\ldots,n_r\in\Bbb Z$.

\spheader 134Yk Tunjukkan bahwa suatu himpunan $E\subseteq\BbbR^r$
terabaikan Lebesgue jika dan hanya jika terdapat suatu barisan
$\sequencen{C_n}$ hiperkubus-hiperkubus dalam $\BbbR^r$ sedemikian sehingga
$E\subseteq\bigcap_{n\in\Bbb N}\bigcup_{k\ge n}C_k$ dan
$\sum_{k=0}^{\infty}(\diam C_k)^r<\infty$, dengan menuliskan $\diam C_k$
untuk diameter $C_k$.
%134Yj, 134F

\spheader 134Yl Tunjukkan bahwa terdapat suatu fungsi kontinu
$f:[0,1]\to[0,1]^2$ sedemikian sehingga $\mu_1f^{-1}[E]=\mu_2E$ untuk
setiap $E\subseteq[0,1]^2$ terukur, dengan menuliskan $\mu_1$, $\mu_2$
untuk ukuran Lebesgue masing-masing pada $\Bbb R$, $\BbbR^2$. \Hint{untuk
setiap $n\in\Bbb N$, nyatakan $[0,1]^2$ sebagai gabungan $4^n$ persegi
tertutup yang panjang sisinya $2^{-n}$; sebut himpunan persegi-persegi ini
$\Cal D_n$. Bangun fungsi-fungsi kontinu $f_n:[0,1]\to[0,1]^2$ dan
keluarga-keluarga $\langle I_D\rangle_{D\in\Cal D_n}$ secara induktif
sedemikian sehingga setiap $I_D$ merupakan interval tertutup yang panjangnya
$4^{-n}$ dan $f_m[I_D]\subseteq D$ setiap kali $D\in\Cal D_n$ dan
$m\ge n$. Induksi akan berjalan lebih lancar jika Anda mengandaikan bahwa
lintasan $f_n$ memasuki setiap persegi dalam $\Cal D_n$ pada salah satu
sudut dan keluar pada sudut yang bersebelahan. Ambil
$f=\lim_{n\to\infty}f_n$. Ini merupakan jenis khusus kurva {\bf Peano} atau
kurva {\bf pengisi ruang}.}
%134H

\spheader 134Ym Tunjukkan bahwa jika $r\le s$, terdapat suatu fungsi kontinu
$f:[0,1]^r\to[0,1]^s$ sedemikian sehingga $\mu_rf^{-1}[E]=\mu_sE$ untuk
setiap $E\subseteq[0,1]^s$ terukur, dengan menuliskan $\mu_r$, $\mu_s$
untuk ukuran Lebesgue masing-masing pada $\BbbR^r$, $\BbbR^s$.
%134Yl, 134H

\spheader 134Yn Tunjukkan bahwa terdapat suatu fungsi kontinu
$f:\Bbb R\to\BbbR^2$ sedemikian sehingga $\mu_1f^{-1}[E]=\mu_2E$ untuk
setiap $E\subseteq\BbbR^2$ terukur, dengan menuliskan $\mu_1$, $\mu_2$
untuk ukuran Lebesgue masing-masing pada $\Bbb R$, $\BbbR^2$.
%134Yl, 134H

\spheader 134Yo Tunjukkan bahwa fungsi $f:[0,1]\to[0,1]^2$ dari 134Yl dapat
dipilih sedemikian sehingga $\mu_2f[E]=\mu_1E$ untuk setiap himpunan terukur
Lebesgue $E\subseteq[0,1]$. \Hint{dengan menggunakan konstruksi yang
disarankan dalam 134Yl, dan menetapkan
$H=f^{-1}[([0,1]\setminus\Bbb Q)^2]$, $f\restr H$ akan merupakan isomorfisme
antara $(H,\mu_{1,H})$ dan $(f[H],\mu_{2,f[H]})$, dengan menuliskan
$\mu_{1,H}$ dan $\mu_{2,f[H]}$ untuk ukuran-ukuran subruang.}
%134Yl, 134H

\spheader 134Yp Tunjukkan bahwa $\Bbb R$ dapat dinyatakan sebagai gabungan
suatu barisan saling lepas $\sequencen{E_n}$ himpunan-himpunan berukuran
berhingga sedemikian sehingga $\mu(I\cap E_n)>0$ untuk setiap interval
terbuka tak-kosong $I\subseteq\Bbb R$ dan setiap $n\in\Bbb N$.
%134J

\spheader 134Yq Tunjukkan bahwa untuk setiap $r\ge 1$, $\BbbR^r$ dapat
dinyatakan sebagai gabungan suatu barisan saling lepas $\sequencen{E_n}$
himpunan-himpunan berukuran berhingga sedemikian sehingga
$\mu(G\cap E_n)>0$ untuk setiap himpunan terbuka tak-kosong
$G\subseteq\BbbR^r$ dan setiap $n\in\Bbb N$.
%134Yp, 134J

\spheader 134Yr Tunjukkan bahwa terdapat suatu barisan saling lepas
$\sequencen{A_n}$ yang terdiri atas subhimpunan-subhimpunan dari $\Bbb R$
sedemikian sehingga
$\mu^*(A_n\cap E)=\mu E$ untuk setiap himpunan terukur $E$ dan setiap
$n\in\Bbb N$.
({\it Catatan\/}: sebenarnya terdapat suatu keluarga saling lepas
$\langle A_t\rangle_{t\in\Bbb R}$ dengan sifat ini, tetapi menurut saya
diperlukan suatu gagasan baru untuk perluasan tersebut. Lihat 419I dalam
Jilid 4.)
%134J

\spheader 134Ys Ulangi 134Yr untuk $\BbbR^r$, dengan $r>1$.
%134Yr, 134J

\spheader 134Yt Uraikan suatu fungsi terukur Borel
$f:[0,1]\to[0,1]$ sedemikian sehingga $f\restr A$ diskontinu pada setiap titik
dari $A$ setiap kali $A\subseteq[0,1]$ merupakan himpunan berukuran luar penuh.
%134J

\spheader 134Yu Misalkan $\sequencen{E_n}$ suatu barisan subhimpunan
terukur tak-terabaikan dari $\BbbR^r$. Tunjukkan bahwa terdapat suatu
himpunan terukur $E\subseteq\BbbR^r$ sedemikian sehingga semua himpunan
$E_n\cap E$, $E_n\setminus E$ tak-terabaikan.
%134J
}%end of exercises

\endnotes{
\Notesheader{134} Ukuran Lebesgue memiliki ragam sifat khusus yang sangat
luas, sesuai dengan kekayaan garis bilangan real beserta struktur aljabar,
topologi, dan keterurutannya. Di sini saya hanya dapat memberi gambaran
sekilas mengenai apa yang mungkin dilakukan.

Ada banyak metode untuk membangun himpunan tak-terukur, dan semuanya
bermakna; metode yang saya berikan dalam 134B mungkin paling mudah diakses,
dan menunjukkan bahwa invariansi terhadap translasi (dengan mengandaikan
aksioma pilihan) merupakan penghalang yang tak-teratasi untuk mengukur setiap
subhimpunan dari $\Bbb R$.

Dalam 134F saya mencantumkan beberapa hubungan dasar antara ukuran dan
topologi ruang Euklides. Hubungan lain terdapat dalam 134Yc, 134Yd, dan 134Yg;
lihat juga 134Xd. Analisis sistematis atas hubungan-hubungan tersebut akan
mengisi sebagian besar Jilid 4.

Himpunan dan fungsi Cantor (134G-134I) membentuk salah satu contoh dasar dalam
teori. Di sini saya menyajikannya sekadar sebagai rancangan menarik dan
sebagai contoh tandingan terhadap suatu dugaan yang wajar. Namun keduanya akan
muncul kembali dalam tiga bab berbeda di Jilid 2 sebagai ilustrasi bagi tiga
fenomena yang sangat berbeda.

Hubungan antara integral Lebesgue dan Riemann jauh lebih dalam daripada yang
ingin saya jelajahi saat ini; fakta bahwa integral Lebesgue memperluas
integral Riemann (134Kb) hanyalah bagian kecil dari ceritanya, dan saya akan
menyesal jika Anda mendapat kesan bahwa integral Lebesgue dengan demikian
menjadikan integral Riemann usang. Tanpa membahas perinciannya di sini, saya
harap 134F dan 134Yg memperjelas bahwa integral Lebesgue dalam arti tertentu
merupakan perluasan kanonik dari integral Riemann. (Setidaknya, saya akan
kembali ke hal ini dalam Bab 43.) Cara lain untuk melihat hal ini adalah
134Yf; integral Lebesgue merupakan integral dasar yang invarian terhadap
translasi pada $\BbbR^r$.
}%end of notes

\discrpage
